\clearpage
\hypertarget{native-476-letter}{}
\begin{center}
\rule{0.55\linewidth}{0.4pt}\\[2pt]
Astronomische Nachrichten. Band XXV. Nr. 582. S. 95. 5. November 1846.\\[2pt]
\rule{0.55\linewidth}{0.4pt}
\end{center}
\begin{flushright}Göttingen, 22. Oktober 1846.\end{flushright}

Gestern Abend habe ich \textsc{Le Verriers} Planeten mit dem Meridiankreis beobachtet. Ich möchte Ihnen das Ergebnis gleich mitteilen:

\begin{center}
1846 Okt. 21.\quad $7^{\rm u}\ 52^{\rm m}\ 13^{\rm s}7$\par
Rektasz. $=327^\circ\ 54'\ 7''8$\qquad Dekl. $-13^\circ\ 32'\ 50''1$
\end{center}

Die Beobachtungen stimmten noch immer alle sehr gut mit einer Kreisbahn überein.
